\section{Заключение}

Настоящата дипломна работа разгледа последователно проектирането и разработването на уеб базирана система за електронна търговия с компютърна и офис техника. В рамките на разработката беше създаден реален Java Spring Boot проект, който покрива основните потребителски и бизнес процеси на един компактен онлайн магазин: регистрация, потвърждение на акаунт, вход, поддръжка на профил, управление на фирми, добавяне на продукти, разглеждане на каталог, работа с количка, финализиране на поръчка и поддържане на история на покупки.

Една от основните цели на труда беше не просто да се реализира работещо приложение, а да се покаже как подобна система може да бъде структурирана по ясен, анализируем и методологично последователен начин. В този смисъл разработката демонстрира важна инженерна стойност. Архитектурата е организирана в отделни слоеве и пакети, домейн моделът е достатъчно ясен и свързан с предметната област, а избраният технологичен стек е едновременно практичен и представителен за съвременна Java уеб разработка.

Проведеният анализ показа, че изборът на Spring Boot, Spring MVC, Thymeleaf, Spring Data JPA и MySQL е напълно оправдан за подобен тип задача. Тези технологии позволяват бързо, но дисциплинирано изграждане на уеб приложение с добра проследимост на логиката от потребителския интерфейс до persistence слоя. Използването на Hibernate Validator, ModelMapper, JavaMail интеграция и \texttt{Pbkdf2PasswordEncoder} допълнително повишава качеството на решението и демонстрира осъзнато използване на наличните средства в екосистемата.

Особено значение в рамките на дипломната работа има фактът, че проектът не беше разгледан само описателно, а критично. Наред с положителните характеристики бяха идентифицирани и конкретни ограничения: опростен security модел, непълна изолация на количката по потребител, силно опростен модел на поръчките, частична локализация и наличие на чувствителни данни в локалната конфигурация. Тези наблюдения са важни, защото показват, че разработката е осмислена не само като демонстрационен резултат, но и като основа за по-нататъшно професионално развитие.

От академична гледна точка могат да бъдат формулирани няколко основни приноса на настоящата работа:

\begin{enumerate}[leftmargin=*]
    \item систематизирано е реално приложение за електронна търговия в последователна многослойна архитектурна рамка;
    \item изведен е ясен модел на предметната област, включващ потребители, фирми, продукти, количка и история на поръчки;
    \item демонстрирано е как чрез Spring технологии може да бъде реализиран пълен потребителски поток в уеб приложение;
    \item извършена е аналитична верификация на функционалността, интерфейса, конфигурацията и operational средата;
    \item формулирани са реалистични насоки за бъдещо развитие към по-сигурна и по-зряла система.
\end{enumerate}

Тези приноси потвърждават, че дипломната работа изпълнява своята двойна задача: от една страна представлява практическа инженерна реализация, а от друга -- завършен технически документ, който обосновава използваните решения. Това е особено важно в контекста на специалност като компютърно и софтуерно инженерство, където способността да се премине от идея към добре структурирана реализация и след това към аргументирана документация е съществен показател за професионална зрялост.

Разработеното приложение не следва да бъде интерпретирано като крайна production-ready търговска платформа. Неговата реална стойност е в това, че представлява работещ, концептуално ясен и надграждаем прототип, върху който могат да бъдат наслагвани следващи итерации: по-силна сигурност, по-богат order management модел, по-чист front-end слой, по-добра конфигурационна дисциплина и автоматизирано тестване. Точно тази отвореност към развитие е един от най-ценните резултати на проекта.

В заключение може да се обобщи, че поставената цел е постигната. Създадена е уеб базирана система за електронна търговия с компютърна и офис техника, която покрива ключов набор от функционалности, реализирана е с подходящи съвременни технологии, описана е в академично издържана форма и е оценена както през призмата на постигнатото, така и през призмата на нейните ограничения. Това прави проекта убедителен завършек на дипломна разработка и стабилна основа за бъдещо инженерно надграждане.
